\documentclass[11pt]{article}
\usepackage[margin=1in]{geometry}
\usepackage{graphicx}
\usepackage{amsmath,amssymb}
\usepackage{caption}
\usepackage{subcaption}
\usepackage{hyperref}
\usepackage{xcolor}
\usepackage{booktabs}
\usepackage{bm}
\usepackage{placeins}  % provides \FloatBarrier
\usepackage{float}     % provides [H] placement

\graphicspath{{../../agent_test_runs/}}

\newcommand{\vA}{v_{\!A}}
\newcommand{\vI}{v_{\!I}}
\newcommand{\vnI}{\mathbf{v}_{n,I}}
\newcommand{\vnA}{\mathbf{v}_{n,A}}
\newcommand{\Ln}{L_n}
\newcommand{\Deff}{D_{\mathrm{eff}}}
\newcommand{\Diso}{D_{\mathrm{iso}}}
\newcommand{\taut}{\tau}
\newcommand{\phin}{\phi_n}
\newcommand{\dx}{\,\mathrm{d}x\,\mathrm{d}y}
\newcommand{\dV}{\,\mathrm{d}V}

\title{Supplementary Information}
\author{}
\date{}

\begin{document}
\maketitle
\tableofcontents
\newpage

%% ======================================================================
\section{Phase-Field Model}
\label{sec:model}
%% ======================================================================

\subsection{Free Energy}

Each cell $n$ in the monolayer is described by a phase field $\phin(\mathbf{x},t)$ that varies smoothly from 1 inside the cell to 0 outside, with an interface of width $\lambda$.  The total free energy is $\mathcal{F} = \mathcal{F}_0 + \mathcal{F}_{\mathrm{int}} + \mathcal{F}_{\mathrm{adh}}$, where the single-cell contribution is
\begin{equation}
\label{eq:F0}
\mathcal{F}_0 = \sum_n \left[\gamma_n \int \left(|\nabla\phin|^2 + \frac{30}{\lambda^2}\phin^2(1-\phin)^2\right)\dV \;+\; \frac{\mu_n}{V_0}\left(V_0 - \int \phin^2\,\dV\right)^{\!2}\right],
\end{equation}
the repulsive cell--cell interaction energy is
\begin{equation}
\label{eq:Fint}
\mathcal{F}_{\mathrm{int}} = \frac{30\kappa}{\lambda^2}\int \sum_{n,\,m\neq n} \phin^2\,\phi_m^2\,\dV,
\end{equation}
and the optional adhesion energy (gradient coupling) is
\begin{equation}
\label{eq:Fadh}
\mathcal{F}_{\mathrm{adh}} = J\sum_{n<m}\int \nabla\phin \cdot \nabla\phi_m\,\dV, \qquad J \ge 0.
\end{equation}
Here $V_0 = \pi R^2$ in 2D and $V_0 = \frac{4}{3}\pi R^3$ in 3D, $\gamma_n$ controls cell stiffness, $\kappa$ controls repulsion strength, $\mu_n$ penalises area/volume changes, and $J$ controls inter-cellular adhesion.

\subsection{Equation of Motion}

The phase field evolves by
\begin{equation}
\label{eq:eom}
\frac{\partial\phin}{\partial t} + \mathbf{v}_n \cdot \nabla\phin = -\frac{1}{2}\frac{\delta\mathcal{F}}{\delta\phin},
\end{equation}
where the factor $\frac{1}{2}$ sets the mobility ($\Gamma = \frac{1}{2}$ in the notation of Palmieri et~al.).  The variational derivative evaluates to
\begin{equation}
\label{eq:varderiv}
\frac{\delta\mathcal{F}}{\delta\phin} = -2\gamma_n\nabla^2\phin + \frac{60\gamma_n}{\lambda^2}\phin(1-\phin)(1-2\phin) - \frac{4\mu_n}{V_0}\!\left(V_0-\!\int\!\phin^2\dV\right)\phin + \frac{60\kappa}{\lambda^2}\phin\!\sum_{m\neq n}\phi_m^2 - J\!\sum_{m\neq n}\nabla^2\phi_m.
\end{equation}

\subsection{Cell Velocity}

The velocity of each cell is
\begin{equation}
\mathbf{v}_n = \vnI + \vnA,
\end{equation}
where the interaction (inactive) velocity is derived from thermodynamic consistency:
\begin{equation}
\label{eq:vI}
\vnI = \frac{60\kappa}{\xi\lambda^2}\int \phin\,(\nabla\phin)\sum_{m\neq n}\phi_m^2\,\dV,
\end{equation}
and the active velocity $\vnA$ has constant magnitude $|\vnA|=\vA$ with orientation governed by run-and-tumble dynamics.  Between reorientation events the cell moves in a fixed direction; the waiting time $t_r$ between successive reorientations is drawn from an exponential distribution,
\begin{equation}
\label{eq:tumble}
P(t_r)\,\mathrm{d}t_r = \frac{1}{\taut}\,e^{-t_r/\taut}\,\mathrm{d}t_r,
\end{equation}
where $\taut$ is the mean persistence time.  Upon reorientation, a new direction is chosen uniformly on the unit circle (2D) or unit sphere (3D).  For an isolated cell, this gives a mean-squared displacement
\begin{equation}
\label{eq:msd_iso}
\langle|\Delta\mathbf{r}|^2\rangle = 2\vA^2\taut\!\left(\frac{t}{\taut} - 1 + e^{-t/\taut}\right),
\end{equation}
which is ballistic at short times and diffusive at long times with effective diffusion coefficient $\Diso = \vA^2\taut/2$.

\subsection{Parameters}

\begin{table}[H]
\centering
\caption{Dimensionless model parameters.}
\label{tab:params}
\begin{tabular}{lccl}
\toprule
Parameter & Symbol & Value & Description \\
\midrule
Interface width     & $\lambda$ & 7          & Phase-field interface thickness \\
Normal stiffness    & $\gamma$  & 1.0        & Gradient energy coefficient \\
Soft stiffness      & $\gamma_{\mathrm{soft}}$ & 0.35 & Cancer/soft cell ($\gamma_{\mathrm{soft}}/\gamma = 0.35$) \\
Repulsion           & $\kappa$  & 10         & Cell--cell overlap penalty \\
Volume constraint   & $\mu$     & 1          & Area/volume conservation strength \\
Target radius       & $R$       & 49         & $V_0 = \pi R^2$ (2D) or $\frac{4}{3}\pi R^3$ (3D) \\
Friction            & $\xi$     & 1500       & Substrate drag coefficient \\
Persistence time    & $\taut$   & $10^4$     & Run-and-tumble reorientation time \\
Active speed        & $\vA$     & $10^{-2}$  & Self-propulsion magnitude \\
Time step           & $\mathrm{d}t$ & 0.02   & Forward Euler step \\
\bottomrule
\end{tabular}
\end{table}

%% ======================================================================
\section{Numerical Methods}
\label{sec:numerics}
%% ======================================================================

\subsection{Spatial Discretisation}

The fields are discretised on a uniform Cartesian grid with spacing $\Delta x = \Delta y = 1$ (2D) or $\Delta x = \Delta y = \Delta z = 1$ (3D).  In 2D, the Laplacian uses the McLellan 9-point isotropic stencil~\cite{McLellan1973}:
\begin{equation}
\nabla^2 f \approx \frac{4(f_N + f_S + f_E + f_W) + (f_{NE}+f_{NW}+f_{SE}+f_{SW}) - 20\,f_C}{6\,h^2},
\end{equation}
which eliminates the $O(h^2)$ grid anisotropy present in the standard 5-point stencil.  In 3D, the analogous 19-point isotropic stencil is used:
\begin{equation}
\nabla^2 f \approx \frac{4(f_{\pm x}+f_{\pm y}+f_{\pm z}) + (f_{\text{12 edges}}) - 36\,f_C}{6\,h^2},
\end{equation}
where the sum includes the 6 face neighbours (weight~4) and 12 edge neighbours (weight~1).  This is the 3D generalisation of the McLellan stencil and eliminates $O(h^2)$ anisotropy in all three coordinate planes.
Gradients use central differences: $\partial_x f \approx (f_{+x}-f_{-x})/(2h)$.

\subsection{Time Integration}

Equation~\eqref{eq:eom} is integrated by forward Euler with $\mathrm{d}t = 0.02$:
\begin{equation}
\phin^{t+1} = \phin^t + \mathrm{d}t\left[-\mathbf{v}_n\cdot\nabla\phin - \tfrac{1}{2}\frac{\delta\mathcal{F}}{\delta\phin}\right].
\end{equation}
The interaction velocity $\vnI$ is computed from the \emph{current} time step's $\phi$ fields via Eq.~\eqref{eq:vI} in a dedicated pre-pass kernel, ensuring no velocity lag in the advection term.

\subsection{Subdomain Decomposition}

Each cell field is stored on a local subdomain (bounding box plus halo) rather than the full global domain.  This reduces per-cell memory from $O(L^d)$ to $O(R^d)$ and enables scaling to large cell counts.  The subdomain tracks the cell centroid and is resized adaptively when the cell approaches the bounding box boundary.

\subsection{GPU Implementation}

The simulation is implemented in CUDA C++.  The main computational kernel (``fused step'') computes, for all cells in parallel: the variational derivative (Laplacian, bulk, constraint, repulsion, adhesion), the advection term, the Euler update, and reduction sums for centroids and perimeters.  Cell--cell interactions are mediated through a global \emph{sum field} $S(\mathbf{x}) = \sum_n \phin^2(\mathbf{x})$, assembled by an atomic scatter kernel before the fused step.  This avoids explicit neighbour lists for the interaction and adhesion terms.  Periodic boundary conditions are applied to both the global domain and individual subdomains.

%% ======================================================================
\section{Extension to Three Dimensions}
\label{sec:3d}
%% ======================================================================

The model extends to 3D with no changes to the free energy functional (Eqs.~\ref{eq:F0}--\ref{eq:Fadh}) or the equation of motion (Eq.~\ref{eq:eom}).  The key differences are:

\begin{enumerate}
\item \textbf{Target volume:} $V_0 = \frac{4}{3}\pi R^3$ replaces $\pi R^2$.
\item \textbf{Laplacian:} The 19-point isotropic stencil (6 face + 12 edge neighbours) replaces the McLellan 9-point stencil, maintaining $O(h^2)$ anisotropy elimination in all three coordinate planes.
\item \textbf{Velocity integral:} Eq.~\eqref{eq:vI} integrates over three spatial dimensions.
\item \textbf{Polarisation:} The active velocity direction $\hat{p}_n$ lives on the unit sphere $S^2$ rather than the circle $S^1$.  Reorientation draws a new direction uniformly on $S^2$.
\item \textbf{Subdomains:} Bounding boxes become 3D ($W\times H\times D$ voxels); the sum-field scatter and fused-step kernels operate on 3D thread grids.
\item \textbf{Memory:} Per-cell storage scales as $O(R^3)$ rather than $O(R^2)$.  For $R=36$ (the 3D default), a subdomain of $\sim 200^3$ voxels requires $\sim$30~MB per cell.
\end{enumerate}

The formulation is otherwise identical.  All validation results in this supplement are for 2D simulations; 3D validation is deferred to the main text.

%% ======================================================================
\section{Validation Against Palmieri et~al.\ (2015)}
\label{sec:palmieri}
%% ======================================================================

We validate the phase-field simulation against Palmieri, Bresler, Wirtz \& Grant, \textit{Sci.\ Rep.}\ \textbf{5}, 11745 (2015), using $N=72$ cells with the parameters in Table~\ref{tab:params}.

\subsection{S4.1.\ Dimensionless Perimeter After Aging (cf.\ Palmieri Fig.~2)}

Monolayers are equilibrated for $8\taut$ at $\vA=0$ from grid-initialised positions.  

\begin{figure}[H]
\centering
\includegraphics[width=\textwidth]{palmieri_fig2_validation.png}
\caption{Monolayer configuration and $\Ln$ distribution after $8\taut$ aging.  Left: $\rho=0.85$; right: $\rho=0.90$.  Top: soft-in-normal ($\gamma_{\mathrm{soft}}=0.35$, green); bottom: all-normal.  The soft cell deforms more at higher packing fraction, consistent with Palmieri Fig.~2.}
\label{fig:aging}
\end{figure}

\subsection{S4.2.\ Tagged-Cell Dynamics (cf.\ Palmieri Fig.~3A)}

Production runs of $48\taut$ with $\vA=0.01$.

\begin{figure}[H]
\centering
\includegraphics[width=\textwidth]{palmieri_fig3a.png}
\caption{Perimeter $\Ln(t)$, centroid speed $|v(t)|$, interaction speed $|\vI(t)|$, and trajectories for the tagged cell.  Top: $\rho=0.85$; bottom: $\rho=0.90$.  Speed bursts above $\vA$ (dashed red) are observed for the soft cell at $\rho=0.90$ but not for the normal cell, consistent with Palmieri Fig.~3A.}
\label{fig:dynamics}
\end{figure}

\subsection{S4.3.\ Velocity Distribution (cf.\ Palmieri Fig.~4)}

\begin{figure}[H]
\centering
\includegraphics[width=\textwidth]{palmieri_fig4.png}
\caption{Scaled velocity distribution $G(v_i)$.  Left: isolated cell (analytical).  Centre: all-normal (Gaussian, $\sigma_G=0.0025$).  Right: soft-in-normal (non-Gaussian tails, Eq.~5 fit in red).  Palmieri reports $\sigma_G\approx 0.0029$ (all-normal) and $0.0028$ (soft-in-normal).}
\label{fig:veldist}
\end{figure}

\subsection{S4.4.\ Effective Diffusion (cf.\ Palmieri Fig.~5)}

\begin{figure}[H]
\centering
\includegraphics[width=0.7\textwidth]{palmieri_fig5.png}
\caption{Instantaneous diffusion $\langle|\Delta\mathbf{r}|^2\rangle/\Delta t$ at $\rho=0.90$ from $48\taut$ runs.  Green: soft cell (soft-in-normal); blue: normal cell (all-normal).  $200\taut$ runs are pending on the cluster for converged $\Deff$ values; Palmieri reports $\Deff^{\mathrm{soft}}/\Deff^{\mathrm{normal}}\approx 1.5$.}
\label{fig:diffusion}
\end{figure}

\subsection{S4.5.\ Summary}

\begin{table}[H]
\centering
\caption{Comparison with Palmieri et~al.\ (2015) at $\rho=0.90$.}
\label{tab:summary}
\begin{tabular}{lcc}
\toprule
Observable & This work ($48\taut$) & Palmieri ($200\taut$) \\
\midrule
$\sigma_G$ (all-normal)           & $0.0025$ & $0.0029$ \\
$\sigma_G$ (soft-in-normal)       & $0.0028$ & $0.0028$ \\
$\langle|v|\rangle/\vA$ (normal)  & $0.31$   & ${\sim}\,0.3$ \\
$\langle|v|\rangle/\vA$ (soft)    & $0.35$   & ${\sim}\,0.35$ \\
Bursts ($|v|>\vA$, soft, $48\taut$) & 4      & 6--8 \\
Bursts ($|v|>\vA$, normal)        & 0        & ${\sim}\,0$ \\
Excess kurtosis (soft)            & $0.51$   & $>0$ \\
Excess kurtosis (normal)          & $0.14$   & ${\approx}\,0$ \\
\bottomrule
\end{tabular}
\end{table}

% PLACEHOLDER: 200τ data
% Once /scratch/ssilber/palmieri_200tau/prod_*/ complete, add:
% - Updated Fig 5 with 200τ curves
% - D_eff(soft)/D_eff(normal) ratio
% - Isolated cell MSD

\begin{thebibliography}{9}
\bibitem{Palmieri2015} B.~Palmieri, Y.~Bresler, D.~Wirtz, and M.~Grant, \textit{Sci.\ Rep.}\ \textbf{5}, 11745 (2015).
\bibitem{Bresler2018} Y.~Bresler, B.~Palmieri, and M.~Grant, arXiv:1807.10318 (2018).
\bibitem{McLellan1973} J.~McLellan, \textit{Proc.\ 7th Annual Princeton Conf.\ Information Sciences and Systems}, p.~247 (1973).
\end{thebibliography}

\end{document}
